\begin{lemma}\label{lemma:E_suffix_preserves_lt}\lean{CollatzMapBasics.E_suffix_preserves_lt}\leanok
Let $k \le j$. If $E_k(n) < E_k(m)$ and the parity bits of $n$ and $m$ agree for steps $i \in \{k, \dots, j-1\}$ (i.e., $X(T^i(n)) = X(T^i(m))$ for $k \le i < j$), then $E_j(n) < E_j(m)$.
\end{lemma}
\begin{proof}\leanok
We proceed by induction on $j$. The base case $j=k$ follows from the hypothesis. For the inductive step $j \to j+1$, we assume $E_j(n) < E_j(m)$ and $X(T^j(n)) = X(T^j(m)) = x$. Then:
\begin{equation}
    E_{j+1}(n) = \frac{3^x}{2} E_j(n) + \frac{x}{2} < \frac{3^x}{2} E_j(m) + \frac{x}{2} = E_{j+1}(m)
\end{equation}
since $3^x/2 > 0$.
\end{proof}

\begin{lemma}\label{lemma:X_of_V_true_false}\lean{CollatzMapBasics.X_of_V_true, CollatzMapBasics.X_of_V_false}\leanok
For any $i < j$, the $i$-th entry of the parity vector $V_j(n)$ is true if and only if $X(T^i(n)) = 1$. It is false if and only if $X(T^i(n)) = 0$.
\end{lemma}
\begin{proof}\leanok
Immediate from the definition of $V_j(n)$.
\end{proof}

\begin{lemma}\label{lemma:E_eq_of_V_prefix_eq}\lean{CollatzMapBasics.E_eq_of_V_prefix_eq}\leanok
If the parity vectors $V_j(m)$ and $V_j(n)$ agree on their first $k$ entries ($k \le j$), then $E_k(m) = E_k(n)$.
\end{lemma}
\begin{proof}\leanok
Since $E_k(n)$ depends only on the parity bits $X(T^0(n)), \dots, X(T^{k-1}(n))$, agreement of the first $k$ entries of the parity vector implies equality of these bits, and thus equality of the correction terms and $E_k$.
\end{proof}

\begin{lemma}\label{lemma:E_elementary_lt}\lean{CollatzMapBasics.E_elementary_lt}\leanok
If $V_j(m)$ and $V_j(n)$ differ by a single elementary swap $01 \to 10$ (i.e., $V_j(m) \prec_{el} V_j(n)$), then $E_j(n) < E_j(m)$.
\end{lemma}
\begin{proof}\leanok
Let the swap occur at positions $k$ and $k+1$. Before position $k$, the parity vectors agree, so $E_k(n) = E_k(m)$. In $V_j(m)$, we have bits $(0, 1)$ at $(k, k+1)$, while in $V_j(n)$ we have $(1, 0)$.
Calculating the two-step increase:
\begin{align*}
    E_{k+2}(m) &= \frac{3^1}{2} \left( \frac{3^0}{2} E_k(m) + \frac{0}{2} \right) + \frac{1}{2} = \frac{3}{4} E_k(m) + \frac{1}{2} \\
    E_{k+2}(n) &= \frac{3^0}{2} \left( \frac{3^1}{2} E_k(n) + \frac{1}{2} \right) + \frac{0}{2} = \frac{3}{4} E_k(n) + \frac{1}{4}
\end{align*}
Thus $E_{k+2}(n) < E_{k+2}(m)$. Since the bits agree for all positions $i > k+1$, Lemma~\ref{lemma:E_suffix_preserves_lt} implies $E_j(n) < E_j(m)$.
\end{proof}

\begin{lemma}\label{lemma:E_pv_transGen_lt}\lean{CollatzMapBasics.E_pv_transGen_lt}\leanok
The correction ratio $E$ is strictly monotonic with respect to the transitive closure of the elementary precedence relation on parity vectors.
\end{lemma}
\begin{proof}\leanok
This follows by induction on the number of elementary swaps, applying Lemma~\ref{lemma:E_elementary_lt} at each step.
\end{proof}

\begin{theorem}[Rozier--Terracol Lemma 2.1]\label{theorem:E_lt_of_V_precedes}\lean{CollatzMapBasics.E_lt_of_V_precedes}\leanok
If $V_j(m)$ strictly precedes $V_j(n)$ in the partial order generated by elementary swaps (i.e., $V_j(m) \prec V_j(n)$), then $E_j(n) < E_j(m)$.
\end{theorem}
\begin{proof}\leanok
This is exactly the statement that the precedence relation $V_j(m) \prec V_j(n)$ implies $E_j(n) < E_j(m)$, which follows from Lemma~\ref{lemma:E_pv_transGen_lt} by identifying the sequence-based $E_j$ with the vector-based calculation.
\end{proof}
